\chapter{複素引数凸関数}
    \section{複素引数凸関数の例}
        \label{複素引数凸関数の例}
        \subsection{\texorpdfstring{$\norm{\sum_{i=1}^N X_i\bm{a}_i + \bm{b}}_2^2\;(\bm{a}_i, \bm{b} \in \complexNumbers^n,X_i \in \complexNumbers^{m\times n})$}{行列X1,X2,...,XNの線形結合×ベクトル+定数のノルム}は\texorpdfstring{$X_1,\dotsc,X_N$}{X1,X2,...,XN}に関して凸である。}
            \begin{proof}
                \quad\par
                $T_X \coloneq (X_1,\dotsc,X_N),T_Y \coloneq (Y_1,\dotsc,Y_N),\lambda \in [0,1]$とし，$f(T_X) \coloneq \norm{\sum_{i=1}^N X_i\bm{a}_i + \bm{b}}_2^2$とする。
                \begin{align*}
                    f((1-\lambda)T_X + \lambda T_Y) &= \norm{\sum_{i=1}^N ((1-\lambda)X_i + \lambda Y_i)\bm{a}_i + \bm{b}}_2^2 \\
                    &= \norm{(1-\lambda)\parens*{\sum_{i=1}^N X_i\bm{a}_i + \bm{b}} + \lambda\parens*{\sum_{i=1}^N Y_i\bm{a}_i + \bm{b}}}_2^2 \\
                    &= (1-\lambda)^2\norm{\sum_{i=1}^N X_i\bm{a}_i + \bm{b}}_2^2 + \lambda^2\norm{\sum_{i=1}^N Y_i\bm{a}_i + \bm{b}}_2^2 \\
                    &\phantom{=} + 2\lambda(1-\lambda)\Re{\parens*{\sum_{i=1}^N X_i\bm{a}_i + \bm{b}}^\mathrm{H}\parens*{\sum_{i=1}^N Y_i\bm{a}_i + \bm{b}}}
                \end{align*}
                よって
                \begin{align*}
                    &\phantom{=} (1-\lambda)f(T_X) + \lambda f(T_Y) - f((1-\lambda)T_X + \lambda T_Y) \\
                    &= \lambda(1-\lambda)\bracks*{\norm{\sum_{i=1}^N X_i\bm{a}_i + \bm{b}}_2^2 + \norm{\sum_{i=1}^N Y_i\bm{a}_i + \bm{b}}_2^2 - 2\Re{\parens*{\sum_{i=1}^N X_i\bm{a}_i + \bm{b}}^\mathrm{H}\parens*{\sum_{i=1}^N Y_i\bm{a}_i + \bm{b}}}} \\
                    &= \lambda(1-\lambda)\norm{\parens*{\sum_{i=1}^N X_i\bm{a}_i + \bm{b}} - \parens*{\sum_{i=1}^N Y_i\bm{a}_i + \bm{b}}}_2^2 \geq 0
                \end{align*}
            \end{proof}
        \subsection{\texorpdfstring{$\norm{A\vx+\bm{b}}_2^2$}{||Ax+b||\string^2}が狭義凸\texorpdfstring{$\iff \HerConj{A}A\succ O$}{⇔ A*Aが正定}}
            \begin{shadebox}
                $A \in \complexNumbers^{m\times n}, \vx \in \complexNumbers^n, \bm{b} \in \complexNumbers^m$とする。
                $f(\vx) \coloneq \norm{A\vx+\bm{b}}_2^2$が$\vx$に関して狭義凸となる必要十分条件は$\HerConj{A}A\succ O$である。
            \end{shadebox}
            $\lambda \in (0,1), \vx,\vy \in \complexNumbers^m, \vx\neq\vy$とする。
            \begin{proof}
                まず次式が成り立つ。
                \begin{align*}
                    f((1-\lambda)\vx + \lambda\vy) &\coloneq \norm{(1-\lambda)(A\vx+\bm{b}) + \lambda(A\vy+\bm{b})}_2^2 \\
                    &= (1-\lambda)^2 f(\vx) + \lambda^2 f(\vy) + 2\lambda(1-\lambda)\Re{\HerConj{(A\vx+\bm{b})}(A\vy+\bm{b})}
                \end{align*}
                これより
                \begin{align*}
                    &\phantom{=} (1-\lambda)f(\vx) + \lambda f(\vy) - f((1-\lambda)\vx + \lambda\vy) \\
                    &= \lambda(1-\lambda)\bracks*{f(\vx) + f(\vy) - 2\Re{\HerConj{(A\vx+\bm{b})}(A\vy+\bm{b})}} \\
                    &= \lambda(1-\lambda)\norm{(A\vx+\bm{b}) - (A\vy+\bm{b})}_2^2 = \lambda(1-\lambda)\norm{A(\vx-\vy)}_2^2 \\
                    &= \lambda(1-\lambda)\HerConj{(\vx-\vy)}\HerConj{A}A(\vx-\vy)
                \end{align*}
                $\lambda \in (0,1),\vx\neq\vy$に対して上式が正となる必要十分条件は$\HerConj{A}A\succ O$である。
            \end{proof}